% Cote 119: offprints, reports, courses, CV and letters (typed and printed part): archive cote 119, site page 167 (PDF page 168).
% typescript. Draft compared: none.
% AI-transcribed by Claude Opus 5.5 (claude-opus-5-5, adaptive thinking, effort high), Message Batches API, batch r119_full (msgbatch_01DwZNm5Y7YsKBWyAN4sFs8w), 2026-11-13; not yet checked by eye.
\scanpage{119/167}
\mspage{\unsure{- 6 -}}

Donnons à $x$ un petit accroissement $h$ ; la différence $F(x+h,y) - F(x,y)$
est égale à $\int_\omega$ le long de n'importe quel chemin d'origine $(x,y)$ et d'extré-%
mité $(x+h,y)$ contenu dans \quad . Si $|h|$ est assez petit on peut prendre pour
tel chemin un segment de droite parallèle à l'axe des $x$.
\[ F(x+h,\,y) - F(x,y) = \int_x^{x+h} P(\quad ,y)\, d \quad \]
Soit \quad $\dfrac{F(x+h,y) - F(x,y)}{h} = \dfrac{1}{h}\displaystyle\int_x^{x+h} P(\quad ,y)\, d \quad$
Lorsque $h$ tend vers $0$, le second membre tend vers $P(x,y)$ en vertu de la
continuité de $P$. On a donc bien $\dfrac{\partial F}{\partial x} = P(x,y)$. On obtient de même $\dfrac{\partial F}{\partial y} = Q(x,y)$
ce qui équivaut d'après (3.2.2) à $\omega = dF$ . C.Q.F.D.
Si $U$ est un disque ouvert nous pouvons renforcer le résultat précédent.
Supposons que l'on sache en effet que $\int_\gamma \omega = 0$ pour les chemins $\gamma$ qui
sont des ``bords'' de rectangles contenus dans $U$ et dont les côtés sont
parallèles aux axes. Un tel rectangle est par exemple pour $\alpha_1 < \alpha_2$ ,
$\beta_1 < \beta_2$ \quad $\{(x,y) \mid \alpha_1 \leq x \leq \alpha_2$ et $\beta_1 \leq y \leq \beta_2\}$. Son bord est le chemin
différentiable par morceau définie par les segments orientés
$[\alpha_1,\beta_1),\ (\alpha_2,\beta_1)],\ [(\alpha_2,\beta_1),(\alpha_2,\beta_2)],\ [(\alpha_2,\beta_2),\ (\alpha_1,\beta_2)],\ [(\alpha_1,\beta_2),\ (\alpha_1,\beta_1)]$.
Notons que d'après (2.1.) il n'est pas nécessaire de préciser ici le paramè-%
trage $\gamma$. Soit $(x_o,y_o)$ le centre du disque $U$ et $(x;y)$ un point quelconque
de $U$. On peut joindre $(x_o,y_o)$ à $(x,y)$ par $\gamma_1$ (segment $[(x_o,y_o),(x,y_o)]$
puis segment $[(x,y_o),\ (x,y)]]$, ou par $\gamma_2$ (analogue par $(x_o,y)$) et d'après
notre hypothèse $\int_{\gamma_1} \omega = \int_{\gamma_2} \omega$ . Désignons cette valeur commune par $F(x;y)$.
La même démonstration que ci-dessus montre que $\dfrac{\partial F}{\partial x} = P$ et $\dfrac{\partial F}{\partial y} = Q$ .
D'où la proposition suivante :
\uline{Proposition 3.4}.

Soit $U$ un disque ouvert. Si l'on a $\int_\gamma \omega = 0$ chaque fois que $\gamma$
est le bord d'un rectangle contenu $U$ et dont les côtés sont parallèles aux
axes, alors $\omega$ admet une primitive dans $U$.

\mspage{\hand{167}}
